\documentclass[11pt,a4paper]{article}
\usepackage[T1]{fontenc}
\usepackage[utf8]{inputenc}
\usepackage{mathpazo}
\usepackage{courier}
\usepackage{amsmath,amssymb,bm}
\usepackage{booktabs,array,tabularx}
\usepackage[margin=2.6cm]{geometry}
\usepackage{microtype}
\usepackage{xcolor}
\usepackage{listings}
\usepackage{upquote} 
\usepackage[authoryear,round]{natbib}
\usepackage[colorlinks=true,linkcolor=blue!50!black,citecolor=blue!50!black,urlcolor=blue!50!black]{hyperref}

\lstset{basicstyle=\ttfamily\small,columns=fullflexible,keepspaces=true,
  frame=single,framerule=0.3pt,rulecolor=\color{black!30},
  backgroundcolor=\color{black!3},xleftmargin=0.6em,xrightmargin=0.6em,
  aboveskip=0.8em,belowskip=0.8em,breaklines=true}
\newcommand{\cmd}[1]{\texttt{#1}}
\newcommand{\yv}{\mathbf y}
\newcommand{\xv}{\mathbf x}
\newcommand{\wv}{\mathbf w}
\newcommand{\E}{\mathrm E}
\setlength{\parskip}{0.35em}
\setlength{\emergencystretch}{3em}

\title{\textbf{hmbivpois}: hidden Markov bivariate Poisson regression\\ for panel counts in Stata\\[0.4em]
\large Technical guide, version 1.0.0}
\author{Paolo Li Donni\\ University of Palermo\\ \href{mailto:paolo.lidonni@unipa.it}{paolo.lidonni@unipa.it}}
\date{November 2026}

\begin{document}
\maketitle

\begin{abstract}
\noindent This guide documents the model, the estimation method and the postestimation tools of
\cmd{hmbivpois}, a Stata command for the hidden Markov bivariate Poisson regression of
\citet{laudicella2022}. The command models two counts observed on the same individuals over time. Given
a latent state that follows a first-order Markov chain, the two counts have Poisson margins whose
means depend on regressors, on the lagged counts and on the state; their dependence within the state
is the form of \citet{lakshminarayana1999} and \citet{famoye2010}, the Frank copula or the normal
copula. Estimation is by the EM algorithm with scaled forward--backward recursions, followed by direct
maximization of the log likelihood. The guide also reports Monte Carlo evidence on the estimator and
on the tests. The command, its help file and an example dataset are available at
\url{https://github.com/plidonni/hmbivpois}.
\end{abstract}

\tableofcontents
\bigskip

%=====================================================================
\section{Introduction}
%=====================================================================

Many outcomes in health economics and elsewhere are pairs of counts recorded repeatedly on the same
individuals: visits to a general practitioner and emergency admissions, doctor visits and hospital
stays, claims of two types. Three sources of association link such counts over time and across the two
outcomes: \emph{state dependence}, the effect of past counts on current ones, both within and between
outcomes; \emph{unobserved heterogeneity}, which may change over time; and \emph{contemporaneous
dependence}, common shocks that hit both counts in the same period.

\citet{laudicella2022} propose a model that separates the three. Unobserved heterogeneity is a latent
state that follows a first-order Markov chain, so that it can change over time; state dependence enters
through the lagged counts; contemporaneous dependence is captured by a bivariate distribution with
Poisson margins. \cmd{hmbivpois} implements this model and extends it in several directions: copula
forms of the dependence, state-specific slopes, transformations of the lagged counts, transition
probabilities that depend on covariates or on time, and a set of postestimation tools.

Section~\ref{sec:model} presents the model, Section~\ref{sec:estimation} the likelihood and its
estimation, Section~\ref{sec:post} inference and postestimation, Section~\ref{sec:command} the
syntax of the command with examples, and Section~\ref{sec:mc} Monte Carlo evidence.

%=====================================================================
\section{The model}\label{sec:model}
%=====================================================================

\subsection{Data and notation}

Individuals are indexed by $i=1,\dots,n$ and periods by $t=0,1,\dots,T_i$. For each individual and
period we observe two counts $\yv_{it}=(y_{1it},y_{2it})$ and regressors $\xv_{1it}$, $\xv_{2it}$ for
the two equations. The first observed period of each individual, $t=0$, is the \emph{presample
period}: it supplies the lagged counts of period~1 and the presample counts $\yv_{i0}$ that enter the
initial probabilities of the states, and it is not an estimation period. The periods of an individual
must be consecutive; panels may be unbalanced.

A latent variable $U_{it}\in\{1,\dots,k\}$ indicates the state of individual $i$ in period $t$. The
model has three blocks: the \emph{measurement model}, the distribution of $\yv_{it}$ given the state;
the \emph{initial probabilities} of the states; and the \emph{transition probabilities}.

\subsection{Measurement model}

Given $U_{it}=c$, the conditional means of the two counts are
\begin{equation}\label{eq:mean}
\log\mu_{jitc}=\alpha_{jc}+\xv_{jit}\bm\beta_j+\gamma_{j1}\,g(y_{1i,t-1})+\gamma_{j2}\,g(y_{2i,t-1}),
\qquad j=1,2,
\end{equation}
where $\alpha_{jc}$ are the state intercepts (support points), $\bm\beta_j$ the slopes of the
regressors and $g(\cdot)$ a transformation of the lagged counts (Section~\ref{sec:lags}). There is no
common constant: the $k$ intercepts of each equation take its place. A single latent chain moves both
intercepts together: the state is a joint characteristic of the individual, such as a health status,
that shifts the demand for both services.

$\gamma_{11}$ and $\gamma_{22}$ measure state dependence \emph{within} each outcome, $\gamma_{12}$ and
$\gamma_{21}$ state dependence \emph{between} outcomes, for example the effect of past primary care on
current emergency care. Because the margins are Poisson in all three forms of dependence (below), the
coefficients keep their usual interpretation as semi-elasticities of $\E(y_{jit}\mid U_{it}=c)$.

\paragraph{State-specific slopes.} With the options \cmd{statevary()} and \cmd{lagvary()}, chosen
slopes in \eqref{eq:mean} become specific to the state, $\bm\beta_{jc}$ or $\gamma_{jlc}$. They are
reported in the equations \textit{depvar}\cmd{\_s1}, \dots, \textit{depvar}\cmd{\_s}$k$.

\subsection{Dependence between the two counts}\label{sec:dep}

The joint probability $p(y_1,y_2\mid\mu_1,\mu_2)$ of the two counts in a state has Poisson margins
with means $\mu_1$ and $\mu_2$ in all three forms. The subscripts $i,t,c$ are dropped in this section.

\paragraph{Famoye form (default).} Following \citet{lakshminarayana1999} and \citet{famoye2010},
\begin{equation}\label{eq:famoye}
p(y_1,y_2)=\frac{\mu_1^{y_1}e^{-\mu_1}}{y_1!}\,\frac{\mu_2^{y_2}e^{-\mu_2}}{y_2!}
\Big[1+\lambda\,\phi_1(y_1)\,\phi_2(y_2)\Big],
\qquad \phi_j(y)=e^{-y}-e^{-q\mu_j},\quad q=1-e^{-1}.
\end{equation}
This is a Sarmanov distribution with mixing functions $\phi_j$. Three properties matter.
\begin{enumerate}
\item \emph{The margins are exactly Poisson.} If $Y\sim\mathrm{Po}(\mu)$, then
$\E(e^{-Y})=e^{-q\mu}$, so $\E\,\phi_j(Y_j)=0$ and $\lambda$ does not affect the margins.
\item \emph{Covariance.} Since $\E\{Y\phi(Y)\}=-q\mu e^{-q\mu}$,
\begin{equation}
\operatorname{Cov}(Y_1,Y_2)=\lambda q^2\mu_1\mu_2\,e^{-q(\mu_1+\mu_2)},\qquad
\rho=\lambda q^2\sqrt{\mu_1\mu_2}\,e^{-q(\mu_1+\mu_2)}.
\end{equation}
The sign of the correlation is that of $\lambda$; its size depends on the means, hence on the
regressors and on the state.
\item \emph{The admissible range of $\lambda$ depends on the means.} With $d_j=e^{-q\mu_j}$, the
functions $\phi_j$ take values in $(-d_j,\,1-d_j]$, and $p\ge0$ for every $(y_1,y_2)$ requires
\begin{equation}\label{eq:range}
-\frac{1}{\max\{d_1d_2,\,(1-d_1)(1-d_2)\}}\;\le\;\lambda\;\le\;
\frac{1}{\max\{(1-d_1)d_2,\,d_1(1-d_2)\}}.
\end{equation}
The interval changes with $(\mu_1,\mu_2)$, and the admissible set is its intersection over all
observations and states. The attainable correlation is small when the means are large: for example,
$|\rho|<0.01$ for $\mu_1=8$, $\mu_2=1$.
\end{enumerate}
\cmd{hmbivpois} keeps $\lambda$ inside the admissible interval during estimation
(Section~\ref{sec:em}) and stores the interval at the estimates in \cmd{e(lambda\_lo)},
\cmd{e(lambda\_hi)} and \cmd{e(lambda\_range)}.

\paragraph{Copula forms.} With \cmd{dependence(frank)} and \cmd{dependence(normal)} the joint
probability is the probability that a copula $C$ gives to the rectangle of the two counts
\citep{cameron2004}:
\begin{equation}\label{eq:copula}
\begin{split}
p(y_1,y_2)={}&C(F_1(y_1),F_2(y_2))-C(F_1(y_1-1),F_2(y_2))\\
&-C(F_1(y_1),F_2(y_2-1))+C(F_1(y_1-1),F_2(y_2-1)),
\end{split}
\end{equation}
where $F_j$ is the Poisson distribution function with mean $\mu_j$. The Frank copula
\citep{frank1979,nelsen2006} is
\begin{equation}
C(u,v;\theta)=-\frac1\theta\log\Big(1+\frac{(e^{-\theta u}-1)(e^{-\theta v}-1)}{e^{-\theta}-1}\Big),
\end{equation}
with $\theta\in\mathbb R$ unrestricted and $\theta=0$ corresponding to independence (reported as
\cmd{theta}). The normal copula is $C(u,v;\rho)=\Phi_2(\Phi^{-1}(u),\Phi^{-1}(v);\rho)$, estimated
on the scale $\operatorname{atanh}\rho$ (reported as \cmd{athrho}; $\rho$ is in \cmd{e(rho)}). The
copula forms allow stronger dependence than \eqref{eq:famoye}, of either sign, at the cost of a more
expensive likelihood.

Two numerical devices keep \eqref{eq:copula} accurate when the cell probability is very small, which
happens often because each observation is evaluated under every state. First, each margin is written
in lower-tail form when $y<\mu$ and in upper-tail form, on the reflected uniform, when $y\ge\mu$; for
both copulas, reflecting one margin only changes the sign of the dependence parameter, so the
rectangle never has its corners close to 1 and its probability is computed without cancellation.
Second, for the normal copula the rectangle probability is computed as a one-dimensional integral
over the margin with the shorter interval,
\begin{equation}
\int_{z_l}^{z_h}\varphi(z)\Big[\Phi\Big(\frac{z'_h-\rho z}{\sqrt{1-\rho^2}}\Big)-\Phi\Big(\frac{z'_l-\rho z}{\sqrt{1-\rho^2}}\Big)\Big]dz,
\end{equation}
by 12-point Gauss--Legendre quadrature; the integrand is positive and smooth, so the result is
accurate in relative terms. Cells with $y_1=y_2=0$ use the bivariate normal distribution function
directly.

\paragraph{One dependence parameter per state.} With \cmd{lambdastate}, $\lambda$ (or $\theta$, or
$\operatorname{atanh}\rho$) is specific to the state.

\paragraph{Local independence.} Given the sequence of states and the lagged counts, the vectors
$\yv_{it}$ are independent over time. Within a period, the two counts are dependent through
$\lambda$ (or the copula parameter). The association between the two counts therefore has three
sources: the lagged counts (state dependence), the persistent latent states (time-varying
heterogeneity) and the within-state dependence (common contemporaneous shocks).

\subsection{Transformations of the lagged counts}\label{sec:lags}

The option \cmd{lagtransform()} sets $g(\cdot)$ in \eqref{eq:mean}, for both lags in both equations
and for the presample counts in the initial probabilities:

\begin{center}
\begin{tabular}{ll}
\toprule
Option & $g(y)$\\
\midrule
\cmd{linear} (default) & $y$\\
\cmd{log1p} & $\log(1+y)$\\
\cmd{logc(\#)} & $\log\max(y,\#)$, $0<\#\le1$\\
\cmd{logzero} & $\log y$ if $y>0$, $0$ otherwise, plus the indicator $\mathbf 1(y=0)$\\
\cmd{indicator} & $\mathbf 1(y>0)$\\
\cmd{topcode(\#)} & $\min(y,\#)$\\
\bottomrule
\end{tabular}
\end{center}

With linear lags and $\gamma_{jj}>0$, the conditional mean grows exponentially in the past count, and
the implied dynamics can be explosive. The log-type transformations make the model a log-linear
autoregression in the counts, whose stability is governed by the matrix of lag coefficients
\citep[see][]{blundell2002}. \cmd{estat stability} (Section~\ref{sec:stability}) diagnoses both
cases.

\subsection{The latent Markov chain}

\paragraph{Initial probabilities.} Following the approach of \citet{heckman1981} and
\citet{wooldridge2005} to initial conditions, the probabilities of the states in the first estimation
period are a multinomial logit on covariates $\xv_{i0}$ (\cmd{initcov()}), on the transformed
presample counts and on a constant, with state~1 as the base:
\begin{equation}
\pi_{ic}=\Pr(U_{i1}=c\mid\xv_{i0},\yv_{i0})=
\frac{\exp(\delta_{0c}+\xv_{i0}\bm\delta_c+g(\yv_{i0})\bm\eta_c)}
{1+\sum_{h=2}^{k}\exp(\delta_{0h}+\xv_{i0}\bm\delta_h+g(\yv_{i0})\bm\eta_h)},\qquad c=2,\dots,k.
\end{equation}
The likelihood is conditional on the presample period. The covariates of \cmd{initcov()} are taken at
the first estimation period. The equations are reported as \cmd{init2}, \dots, \cmd{init}$k$.

\paragraph{Transition probabilities.} Each row of the transition matrix is a multinomial logit with
staying in the state as the base:
\begin{equation}\label{eq:trans}
\pi_{d\mid c,it}=\Pr(U_{it}=d\mid U_{i,t-1}=c)=
\frac{\exp(\tau_{cd}+\wv_{it}\bm\psi_{cd})}{1+\sum_{h\ne c}\exp(\tau_{ch}+\wv_{it}\bm\psi_{ch})},
\qquad d\ne c,
\end{equation}
reported in the equations \cmd{tr}$c$\cmd{\_}$d$. Four specifications are available:
\begin{itemize}
\item \emph{homogeneous} (default): no $\wv_{it}$, $k(k-1)$ parameters, as in \citet{laudicella2022};
\item \emph{covariates}, \cmd{xtrans(}\textit{varlist}\cmd{)}: $\wv_{it}$ measured at the period of
arrival $t$;
\item \emph{time-varying}, \cmd{tvtrans}: period dummies in place of the constant, that is, a free
transition matrix for each period of arrival; it can be combined with \cmd{xtrans()};
\item \emph{no transitions}, \cmd{notransition}: the state of each individual is constant, which gives
a latent class (finite mixture) model with the same measurement model.
\end{itemize}

\subsection{Number of parameters}

With $p_j$ common slopes, $s_j$ state-specific slopes, $m$ covariates in \cmd{initcov()} and $r$
covariates of the transitions, the model has
\[
\underbrace{\textstyle\sum_j(p_j+k+k s_j)}_{\text{measurement}}+\underbrace{n_\lambda}_{1\text{ or }k}
+\underbrace{(k-1)(m+L+1)}_{\text{initial}}+\underbrace{k(k-1)(1+r)}_{\text{transitions}}
\]
free parameters, where $L$ is the number of features of the transformed presample counts; with \cmd{tvtrans} the constant of each transition logit is replaced by one dummy per
period of arrival.

%=====================================================================
\section{Likelihood and estimation}\label{sec:estimation}
%=====================================================================

\subsection{Likelihood}

Let $f_{itc}=p(\yv_{it}\mid U_{it}=c,\xv_{it},\yv_{i,t-1})$. The likelihood of individual $i$ sums
over all the sequences of states,
\begin{equation}
L_i=\sum_{c_1=1}^k\cdots\sum_{c_{T_i}=1}^k\pi_{ic_1}f_{i1c_1}\prod_{t=2}^{T_i}\pi_{c_t\mid c_{t-1},it}\,f_{itc_t},
\qquad \ell(\bm\vartheta)=\sum_{i=1}^n\log L_i,
\end{equation}
and is computed by the forward recursion \citep{baum1970,zucchini2016}
\begin{equation}
a_{i1}(c)=\pi_{ic}f_{i1c},\qquad
a_{it}(d)=\Big[\sum_{c}a_{i,t-1}(c)\,\pi_{d\mid c,it}\Big]f_{itd},\qquad
L_i=\sum_c a_{iT_i}(c).
\end{equation}
The recursion is scaled: at each $t$ the vector $a_{it}$ is normalized to sum to one and the log of
the scale factor is accumulated, so that long panels do not underflow. Individuals are sorted by
decreasing $T_i$, so that those observed at period $t$ are the first rows of the data; the recursions
are then vectorized over individuals, with no loop over $i$, and unbalanced panels need no special
treatment.

\subsection{The EM algorithm}\label{sec:em}

Estimation follows the EM algorithm \citep{dempster1977} as adapted to latent Markov models by
\citet{bartolucci2013}. The complete-data log likelihood separates into three blocks: the measurement
model, the initial probabilities and the transitions.

\paragraph{E-step.} The forward--backward recursions give the posterior probabilities of the states,
$\tilde w_{itc}=\Pr(U_{it}=c\mid\yv_i)$, and the joint posteriors of consecutive states,
$\tilde z_{it}(c,d)=\Pr(U_{i,t-1}=c,U_{it}=d\mid\yv_i)$.

\paragraph{M-step.} Each block is updated given the posteriors.
\begin{itemize}
\item \emph{Measurement model:} the bivariate regression with state intercepts, with weights
$\tilde w_{itc}$, is updated by Newton steps with step halving until the expected complete log
likelihood does not decrease. With the Famoye form the Hessian is analytic and two steps are taken
per iteration; with the copulas minus the outer product of the scores replaces the second
derivatives and one step is taken, because each evaluation costs more. With the Famoye form, every
trial point is made admissible by projecting $\lambda$ onto the interval \eqref{eq:range} at the trial
means; a $\lambda$ on a bound with the gradient pointing outwards is left out of the Newton system
(active set), so that the other parameters keep moving.
\item \emph{Initial probabilities:} a fractional multinomial logit with responses $\tilde w_{i1c}$,
fitted by Newton--Raphson.
\item \emph{Transitions:} with the homogeneous chain the update has the closed form
$\hat\pi_{d\mid c}\propto\sum_{i,t}\tilde z_{it}(c,d)$; with \cmd{xtrans()} or \cmd{tvtrans}, one
weighted multinomial logit per origin state $c$, with responses $\tilde z_{it}(c,\cdot)$.
\end{itemize}
The algorithm stops when the relative change of the log likelihood falls below
\cmd{emtolerance()} (default $10^{-8}$) or after \cmd{emiterate()} iterations (default 500).

\paragraph{Direct maximization.} If the EM has not converged, or if \cmd{technique()} is specified,
the marginal log likelihood is maximized directly with \cmd{optimize()} (BFGS by default). The
analytic score is computed by the Fisher identity: the score of $\ell$ equals the expected
complete-data score under the posteriors, at the cost of one forward--backward pass.

\paragraph{Starting values.} By default the slopes start from univariate Poisson regressions of the
two counts, the state intercepts from the Poisson constant plus offsets spread evenly between $-0.8$
and $0.8$, the dependence parameter from zero, the initial logit from zero and the chain from a
persistent matrix with staying probability 0.9. Starting values can be given with \cmd{from()}. As
with any finite mixture, the likelihood can have local maxima; trying several starting values is
advisable in applications.

\paragraph{Labelling of the states.} After estimation the states are relabelled in increasing order of
the sample mean of the fitted $\mu_{1c}$ (\cmd{sortby(eq1)}, the default) or $\mu_{2c}$
(\cmd{sortby(eq2)}); with common slopes this is the order of the state intercepts. This makes the
output stable across runs and comparable across values of $k$.

\subsection{Variance}

The default variance is the inverse of the observed information, computed by central differences of
the analytic score \citep{mclachlan2000}. A Famoye $\lambda$ that lies on a bound of its admissible
interval is held fixed: its rows and columns of $\mathbf V$ are zero, and \cmd{e(n\_bound)} counts such
parameters. \cmd{vce(robust)} gives the sandwich estimator clustered on the individuals, built from the
individual contributions to the score.

%=====================================================================
\section{Inference and postestimation}\label{sec:post}
%=====================================================================

\subsection{Number of states}

\cmd{estat ic} reports the log likelihood, the number of parameters, AIC and
$\mathrm{BIC}=-2\ell+\#\text{par}\cdot\log n$ with $n$ the number of \emph{individuals}, the standard
choice in latent Markov models \citep{bartolucci2013}; \cmd{n()} changes it. It also reports the
relative entropy of the classification,
\[
1-\frac{\sum_{i,t,c}-\tilde w_{itc}\log\tilde w_{itc}}{N\log k},
\]
which equals one when the states are known with certainty ($N$ is the number of estimation periods).

\subsection{Transition probabilities}

\cmd{estat transition} reports the transition matrix, the expected duration of a stay in each state,
$1/(1-\pi_{c\mid c})$, and the mean initial probabilities, with delta-method standard errors; the
confidence intervals of the probabilities are computed on the logit scale. With heterogeneous
transitions the reported matrix is the mean over the transition rows, or over those into one period
with \cmd{period()}.

\subsection{Test of no transitions}

Under $H_0$: $\pi_{d\mid c}=0$ for all $d\ne c$, the states are constant and the model is a latent
class model. $H_0$ lies on the boundary of the parameter space, so the likelihood-ratio statistic does
not have the usual $\chi^2$ distribution. Following \citet{bartolucci2006}, its null distribution is a
chi-bar-squared, a mixture
\[
\Pr(LR\le x)=\sum_{j=0}^{k(k-1)}\omega_j\,\Pr(\chi^2_j\le x),
\]
whose weight $\omega_j$ is the probability that the projection of a $N(\mathbf 0,\mathbf I^{-1})$
vector onto the nonnegative orthant has $j$ positive components, where $\mathbf I$ is the information
on the off-diagonal transition probabilities net of the other parameters, at the restricted estimates.
\cmd{estat lrtrans} fits the restricted model (twice, from the default starting values and from the
unrestricted estimates, keeping the larger likelihood), estimates $\mathbf I$ from the outer product of
the individual scores, and estimates the weights by \cmd{reps()} simulated projections (default
10\,000), each solved as a nonnegative least squares problem \citep{lawson1974}. It also reports the
$p$-value from the $\chi^2_{k(k-1)}$ distribution, which is conservative.

\subsection{Predictions}

\cmd{predict} computes, for the estimation sample, statistics that average over the states with
weights given by one of four sets of state probabilities:

\begin{center}
\begin{tabularx}{\textwidth}{lX}
\toprule
Weights & Probability of $U_{it}=c$ given\\
\midrule
posterior (default) & all the outcomes of the individual, $\tilde w_{itc}$\\
\cmd{filtered} & the outcomes up to period $t$\\
\cmd{onestep} & the outcomes up to period $t-1$\\
\cmd{prior} & the covariates and the presample counts only\\
\cmd{instate(\#)} & the statistic in state \# \\
\bottomrule
\end{tabularx}
\end{center}

The statistics are the expected counts \cmd{mu1}, \cmd{mu2}; the state probabilities \cmd{postprob};
the most likely state of each period \cmd{state}; the most likely sequence of states of each
individual \cmd{viterbi}, by the Viterbi algorithm in logs \citep{viterbi1967}; the within-state
correlation \cmd{rho}; and the probability of a pair of counts \cmd{pr(}$a,b$\cmd{)}. For the
copulas, $\rho$ is computed from
$\operatorname{Cov}(Y_1,Y_2)=\sum_{a,b\ge0}[C(F_1(a),F_2(b))-F_1(a)F_2(b)]$, summed until the upper
tails of both margins are negligible. The one-step weights give a check on the fit: the sum over the
sample of $\log\cmd{pr}(y_{1it},y_{2it})$ with \cmd{onestep} equals the log likelihood.

\subsection{Average marginal effects}

\cmd{estat ame} reports average marginal effects on $\E(y_1)$ and $\E(y_2)$ with the state
probabilities held fixed:
\[
\mathrm{AME}=\frac1N\sum_{i,t}\sum_{c}\Pr(U_{it}=c\mid\cdot)\,
\frac{\mu_{jitc}(A)-\mu_{jitc}(B)}{\Delta},
\]
where $A$ and $B$ are two counterfactual designs: the regressor shifted by $\pm h$ for continuous
variables ($\Delta=2h$), a level and the base level for factor variables ($\Delta=1$), one more count
of a lagged outcome ($\Delta=1$). The weights are those of \cmd{predict}. Standard errors are by the
delta method. Variables coded 0/1 without \cmd{i.} are treated as continuous. \cmd{margins} is not
supported.

\subsection{Stability of the dynamics}\label{sec:stability}

\cmd{estat stability} checks whether the dynamics of the lagged counts can explode. It reports the lag
coefficients by state and:
\begin{itemize}
\item with linear lags, the index $\gamma_{jj}\exp(a)$ of the univariate map $y\mapsto\exp(a+\gamma_{jj}y)$,
where $a$ is a high quantile of the rest of the linear predictor; the map has a fixed point only if the
index does not exceed $e^{-1}$;
\item with log-type lags, the spectral radius of the $2\times2$ matrix of the coefficients of
$g(y_1)$ and $g(y_2)$, which must be below one;
\item the share of deterministic paths, from every observation and in every state, whose expected count
exceeds \cmd{limit()} (default 1000) within \cmd{horizon()} periods (default 50);
\item the same share for \cmd{nsim()} simulated paths per individual of the estimated model, with the
state drawn from its posterior probabilities and the estimated chain.
\end{itemize}

%=====================================================================
\section{The command}\label{sec:command}
%=====================================================================

\subsection{Syntax}

\begin{lstlisting}
hmbivpois (depvar1 = [indepvars1] [, eqopts]) (depvar2 = [indepvars2] [, eqopts])
          [if] [in], id(varname) time(varname) k(#) [options]
\end{lstlisting}

The data are in long form, one observation per individual and period. \textit{indepvars},
\cmd{statevary()} and \cmd{xtrans()} may contain factor variables. The main options are:

\begin{center}
\small
\begin{tabularx}{\textwidth}{lX}
\toprule
Option & Description\\
\midrule
\multicolumn{2}{l}{\emph{Equation options}}\\
\cmd{statevary(}\textit{varlist}\cmd{)} & regressors with state-specific coefficients\\
\cmd{lagvary(}\textit{depvars}\cmd{)} & lagged outcomes with state-specific coefficients\\
\midrule
\multicolumn{2}{l}{\emph{Model}}\\
\cmd{id()}, \cmd{time()}, \cmd{k(\#)} & individual, integer period, number of states (required)\\
\cmd{initcov(}\textit{varlist}\cmd{)} & covariates of the initial probabilities\\
\cmd{dependence(}\textit{form}\cmd{)} & \cmd{famoye} (default), \cmd{frank} or \cmd{normal}\\
\cmd{lambdastate} & one dependence parameter per state\\
\cmd{lagtransform(}\textit{g}\cmd{)} & transformation of the lagged counts (Section~\ref{sec:lags})\\
\midrule
\multicolumn{2}{l}{\emph{Transitions}}\\
\cmd{xtrans(}\textit{varlist}\cmd{)} & covariates of the transitions, at the period of arrival\\
\cmd{tvtrans} & a free transition matrix for each period of arrival\\
\cmd{notransition} & no transitions (latent class model)\\
\midrule
\multicolumn{2}{l}{\emph{Estimation and reporting}}\\
\cmd{emiterate(\#)}, \cmd{emtolerance(\#)} & EM iterations (500) and tolerance ($10^{-8}$)\\
\cmd{emonly}, \cmd{technique()}, \cmd{iterate(\#)} & control of the direct maximization\\
\cmd{from(}\textit{matname}\cmd{)} & starting values, in the order of \cmd{e(b)}\\
\cmd{sortby(eq1|eq2|none)} & labelling of the states\\
\cmd{vce(oim|robust)}, \cmd{level(\#)}, \cmd{nolog} & variance and display\\
\bottomrule
\end{tabularx}
\end{center}

Postestimation:
\begin{lstlisting}
predict [type] newvar | stub* [if] [in] [, mu1 mu2 postprob state viterbi rho pr(a,b)
        filtered onestep prior instate(#)]
estat ic [, n(#)]
estat transition [, period(#) level(#)]
estat lrtrans [, reps(#) seed(#)]
estat stability [, horizon(#) limit(#) nsim(#) seed(#)]
estat ame [, prior onestep filtered instate(#) level(#)]
\end{lstlisting}

The stored results are listed in \cmd{help hmbivpois}.

\subsection{Installation}

\begin{lstlisting}
net install hmbivpois, from("https://raw.githubusercontent.com/plidonni/hmbivpois/main/") replace
net get hmbivpois, from("https://raw.githubusercontent.com/plidonni/hmbivpois/main/") replace
\end{lstlisting}
The first command installs the command and its help file; the second downloads into the current
folder the example dataset \cmd{hmbivpois\_example.dta} and the do-file \cmd{hmbivpois\_examples.do},
which runs all the examples below. Stata 16 or later is required.

\subsection{Examples}

The example dataset is simulated from the model with $n=1000$ individuals, a presample period and
$T=6$ estimation periods, $k=3$ states, the Famoye form with $\lambda=-0.5$ and homogeneous
transitions. The true values are listed in Table~\ref{tab:true}; the variable \cmd{state} holds the
true latent state.

\begin{table}[h]
\centering\small
\caption{True parameters of the example dataset}\label{tab:true}
\begin{tabular}{ll}
\toprule
Block & Values\\
\midrule
Equation $y_1$ & $x_{11}=0.4$, $x_{12}=-0.3$, $\gamma_{11}=0.03$, $\gamma_{12}=-0.05$, $\alpha_{1}=(-0.3,\,0.7,\,1.5)$\\
Equation $y_2$ & $x_{21}=0.3$, $x_{22}=-0.4$, $\gamma_{21}=-0.10$, $\gamma_{22}=0.10$, $\alpha_{2}=(-2.5,\,-1.5,\,-0.6)$\\
Dependence & $\lambda=-0.5$\\
Initial (states 2, 3) & $z_1=(0.5,1.0)$, $y_{1,0}=(0.15,0.30)$, $y_{2,0}=(0.30,0.50)$, constant $=(-0.5,-2.0)$\\
Transitions & $\Pi=\begin{pmatrix}0.85&0.12&0.03\\0.08&0.84&0.08\\0.03&0.12&0.85\end{pmatrix}$\\
\bottomrule
\end{tabular}
\end{table}

\medskip\noindent\textbf{Base model.}\par\nopagebreak
\begin{lstlisting}
use hmbivpois_example, clear
hmbivpois (y1 = x11 x12) (y2 = x21 x22), id(id) time(time) k(3) initcov(z1)
\end{lstlisting}
The output reports, for each equation, the slopes, the lag coefficients (\cmd{lag\_y1},
\cmd{lag\_y2}) and the state intercepts (\cmd{state1}, \dots), followed by $\lambda$, the initial logit
(\cmd{init2}, \cmd{init3}) and the transition logits (\cmd{tr1\_2}, \dots).

\medskip\noindent\textbf{Postestimation.}\par\nopagebreak
\begin{lstlisting}
estat ic
estat transition
estat ame
predict double pp*, postprob
predict s_hat, viterbi
tabulate state s_hat if time > 0
estat lrtrans
\end{lstlisting}
With simulated data, the cross-tabulation of the true state and of the Viterbi path shows how well the
states are recovered.

\medskip\noindent\textbf{Choice of the number of states.}\par\nopagebreak
\begin{lstlisting}
forvalues k = 1/4 {
    quietly hmbivpois (y1 = x11 x12) (y2 = x21 x22), id(id) time(time) k(`k') initcov(z1) nolog
    estat ic
}
\end{lstlisting}

\medskip\noindent\textbf{Extensions.}\par\nopagebreak
\begin{lstlisting}
hmbivpois (y1 = x11 x12, statevary(x11)) (y2 = x21 x22), id(id) time(time) ///
    k(3) initcov(z1) dependence(frank) xtrans(z1) lagtransform(log1p)
estat stability
\end{lstlisting}
This specification uses the Frank copula, lets the slope of $x_{11}$ differ across states, makes the
transitions depend on $z_1$ and transforms the lags by $\log(1+y)$. It is misspecified for the example
data, which were generated without these features, and serves only to illustrate the syntax.

\paragraph{Practical advice.}
\begin{itemize}
\item Fit the model for increasing $k$ and compare the BIC; check the entropy and the size of the
smallest state.
\item Re-estimate with different starting values (\cmd{from()}) to guard against local maxima,
especially for $k\ge3$.
\item With the Famoye form, a $\lambda$ on a bound of its admissible interval (\cmd{e(n\_bound)}$>0$)
signals that the data call for stronger dependence than the form allows: try a copula.
\item With linear lags and large counts, run \cmd{estat stability}, or use a log-type
\cmd{lagtransform()}.
\end{itemize}

%=====================================================================
\section{Monte Carlo evidence}\label{sec:mc}
%=====================================================================

The command was tested with simulation experiments, each with $n=2000$ individuals, a presample period,
$T=8$ estimation periods and $k=3$ states, using the design of Table~\ref{tab:true} for the
measurement model, the initial probabilities and the chain, with the changes listed in
Table~\ref{tab:mc}. For the estimators, the table reports the ratio between the Monte Carlo standard
deviation of the estimates and the mean estimated standard error (one if the standard errors are
correct), and the coverage of the nominal 95\% confidence intervals, as ranges over all the
parameters.

\begin{table}[h]
\centering\small
\caption{Monte Carlo results ($n=2000$, $T=8$, $k=3$)}\label{tab:mc}
\begin{tabularx}{\textwidth}{>{\raggedright\arraybackslash}p{4.3cm}cXX}
\toprule
Design & Reps & MC s.d./mean s.e.: mean [min, max] & Coverage of 95\% CI [min, max]\\
\midrule
Frank copula, $\theta=3$ & 100 & 0.998 [0.879, 1.095] & [0.920, 0.990]\\
Normal copula, $\rho=0.5$ & 200 & 1.008 [0.935, 1.100] & [0.925, 0.985]\\
Transition probabilities (\cmd{estat transition}), Famoye, $\lambda=-0.5$ & 100 & 0.947 [0.833, 1.031] & [0.940, 0.980]\\
Full model: Frank $\theta=3$, slope of $x_{11}$ by state, transitions on a covariate & 100 & 1.007 [0.904, 1.101] & [0.910, 0.980]\\
\bottomrule
\end{tabularx}
\end{table}

\paragraph{Test of no transitions.} With data generated without transitions (Famoye, $\lambda=-0.5$,
100 replications), the rejection rates of \cmd{estat lrtrans} at the 10\%, 5\% and 1\% levels were
0.100, 0.050 and 0.010 with the chi-bar-squared distribution, against 0.010, 0.010 and 0.000 with the
$\chi^2_6$ distribution, which is markedly conservative. The mean LR statistic was 2.20, close to the
mean of the estimated chi-bar-squared distribution (2.56) and far from that of $\chi^2_6$.

\paragraph{Other checks.} With $\theta=0$ the Frank copula and with $\lambda=0$ the Famoye form give the
same log likelihood, as they must under independence. With the Frank copula, the likelihood-ratio test
of a common $\theta$ against one per state did not reject when the true $\theta$ was common
($\chi^2_2=1.58$, $p=0.45$).

\paragraph{Caveat.} In the full-model design the largest bias relative to its Monte Carlo standard error
(4.6) concerns the constant of the transition from state 1 to state 3, the rarest transition of the
design (probability 0.03); its bias, $-0.098$, is small relative to its standard error (0.21). With
100 replications such a result is compatible with finite-sample bias of a parameter identified by few
transitions, but it suggests caution when interpreting rare transitions in small samples.

%=====================================================================
\section*{Citation}
%=====================================================================
If you use \cmd{hmbivpois}, please cite \citet{laudicella2022} and

\smallskip
\noindent Li Donni, P. 2026. hmbivpois: Stata module to fit hidden Markov bivariate Poisson regression for
panel counts. \url{https://github.com/plidonni/hmbivpois}.

\begin{thebibliography}{99}
\bibitem[Bartolucci(2006)]{bartolucci2006}
Bartolucci, F. 2006. Likelihood inference for a class of latent Markov models under linear hypotheses on
the transition probabilities. \textit{Journal of the Royal Statistical Society, Series B} 68: 155--178.

\bibitem[Bartolucci et~al.(2013)]{bartolucci2013}
Bartolucci, F., A. Farcomeni, and F. Pennoni. 2013. \textit{Latent Markov Models for Longitudinal Data}.
Boca Raton, FL: Chapman \& Hall/CRC.

\bibitem[Baum et~al.(1970)]{baum1970}
Baum, L. E., T. Petrie, G. Soules, and N. Weiss. 1970. A maximization technique occurring in the
statistical analysis of probabilistic functions of Markov chains. \textit{Annals of Mathematical
Statistics} 41: 164--171.

\bibitem[Blundell et~al.(2002)]{blundell2002}
Blundell, R., R. Griffith, and F. Windmeijer. 2002. Individual effects and dynamics in count data
models. \textit{Journal of Econometrics} 108: 113--131.

\bibitem[Cameron et~al.(2004)]{cameron2004}
Cameron, A. C., T. Li, P. K. Trivedi, and D. M. Zimmer. 2004. Modelling the differences in counted
outcomes using bivariate copula models with application to mismeasured counts. \textit{Econometrics
Journal} 7: 566--584.

\bibitem[Dempster et~al.(1977)]{dempster1977}
Dempster, A. P., N. M. Laird, and D. B. Rubin. 1977. Maximum likelihood from incomplete data via the EM
algorithm. \textit{Journal of the Royal Statistical Society, Series B} 39: 1--38.

\bibitem[Famoye(2010)]{famoye2010}
Famoye, F. 2010. On the bivariate negative binomial regression model. \textit{Journal of Applied
Statistics} 37: 969--981.

\bibitem[Frank(1979)]{frank1979}
Frank, M. J. 1979. On the simultaneous associativity of $F(x,y)$ and $x+y-F(x,y)$.
\textit{Aequationes Mathematicae} 19: 194--226.

\bibitem[Heckman(1981)]{heckman1981}
Heckman, J. J. 1981. The incidental parameters problem and the problem of initial conditions in
estimating a discrete time--discrete data stochastic process. In \textit{Structural Analysis of
Discrete Data with Econometric Applications}, ed. C. F. Manski and D. McFadden, 179--195. Cambridge, MA:
MIT Press.

\bibitem[Lakshminarayana et~al.(1999)]{lakshminarayana1999}
Lakshminarayana, J., S. N. N. Pandit, and K. Srinivasa Rao. 1999. On a bivariate Poisson distribution.
\textit{Communications in Statistics -- Theory and Methods} 28: 267--276.

\bibitem[Laudicella and Li Donni(2022)]{laudicella2022}
Laudicella, M., and P. Li Donni. 2022. The dynamic interdependence in the demand of primary and emergency
secondary care: A hidden Markov approach. \textit{Journal of Applied Econometrics} 37: 521--536.

\bibitem[Lawson and Hanson(1974)]{lawson1974}
Lawson, C. L., and R. J. Hanson. 1974. \textit{Solving Least Squares Problems}. Englewood Cliffs, NJ:
Prentice-Hall.

\bibitem[McLachlan and Peel(2000)]{mclachlan2000}
McLachlan, G., and D. Peel. 2000. \textit{Finite Mixture Models}. New York: Wiley.

\bibitem[Nelsen(2006)]{nelsen2006}
Nelsen, R. B. 2006. \textit{An Introduction to Copulas}. 2nd ed. New York: Springer.

\bibitem[Viterbi(1967)]{viterbi1967}
Viterbi, A. J. 1967. Error bounds for convolutional codes and an asymptotically optimum decoding
algorithm. \textit{IEEE Transactions on Information Theory} 13: 260--269.

\bibitem[Wooldridge(2005)]{wooldridge2005}
Wooldridge, J. M. 2005. Simple solutions to the initial conditions problem in dynamic, nonlinear panel
data models with unobserved heterogeneity. \textit{Journal of Applied Econometrics} 20: 39--54.

\bibitem[Zucchini et~al.(2016)]{zucchini2016}
Zucchini, W., I. L. MacDonald, and R. Langrock. 2016. \textit{Hidden Markov Models for Time Series: An
Introduction Using R}. 2nd ed. Boca Raton, FL: Chapman \& Hall/CRC.
\end{thebibliography}

\end{document}
